\section{A breathable Moon as a modelling regime}

The models take Terluna's design scenario as fixed conditions. The air holds 121,590~Pa at the water level,
with oxygen at Earth's partial pressure of 21,227~Pa (17.5\% of the air), argon in Earth's ratio to nitrogen,
400~ppm of carbon dioxide and nitrogen for the balance; the land stands above the water level, so the
global-mean surface pressure is 1.16~atm and ranges from 1.02 to 1.22~atm.\claim{S05} Seas fill the lunar
basins to 28\% of the surface and rain-fed lakes stand at their own levels, so 759 of the GCM's 2,048
cells, 40\% of its surface, hold water.\claim{S06} Sunlight is the TSIS-1 solar spectrum
\autocite{coddington2021tsis} passed through a titania shield in space, 1,235~W\,m$^{-2}$ at the top of the
air, dimmed evenly at every wavelength to 1,198~W\,m$^{-2}$ in the design case; the shield removes the
ultraviolet that would make ozone, so the air carries none.\claim{S07} The land has a placeholder albedo of
0.2 and the seas a 50~m slab with sea ice. Every climate number below holds given these choices; they are a
design, and other designs would move the numbers.

\begin{table}[t]
\caption{The ratios that separate the Open Moon from Earth, and the model parts each strains. Lunar values
are for the 1.2~atm design case; scale heights use the surface temperature and the air's mean molar mass;
horizon air masses are Chapman's grazing factor $\sqrt{\pi R/2H}$ for an isothermal air, which PlaSim's code
sets at 35 for Earth. The convective grid is a nominal scale from the dry dynamics' similarity.}
\label{tab:regime}
\small
\begin{tabularx}{\textwidth}{@{}lrrL@{}}
\toprule
Quantity & Earth & Open Moon & What it strains \\
\midrule
Surface gravity (m\,s$^{-2}$) & 9.81 & 1.62 & Dynamics; fall speeds; sea roughness; wave dispersion; heights in pressure terms \\
Surface pressure at the water level (bar) & 1.01 & 1.22 & RRTMG's tables, whose highest pressure is 1.05~bar \\
Air over each square metre (t) & 10.3 & 74.9 & Broadband absorber paths fitted on Earth's \\
Scale height at the surface (km) & 8.4 & 52 & Every height-based formula \\
Scale height over radius & 0.0013 & 0.030 & Flat, sunlit-only columns; the thin shell \\
Horizon air mass of well-mixed air & 34.5 & 7.2 & Earth's slant factor, fixed in PlaSim \\
Solar day (Earth days) & 1 & 29.53 & Formulas written per day; the soil's thermal wave; the model calendar \\
Sidereal rotation (Earth days) & 0.997 & 27.32 & ExoPlaSim's Sun clock; the Coriolis parameter \\
Tilt of the equator (\textdegree) & 23.44 & 1.54 & ExoPlaSim's orbit default \\
Height of the 0\,\textdegree C level (km) & 4.5--5 & 25 & Ground-strike rules; the charging zone \\
Tops of deep storm cloud (km) & 15--18 & 55--80 & Model tops; the levels of cold cloud \\
Fall speed of rain, relative & 1 & 0.34 & Evaporation of falling rain \\
Grid for deep convection (km) & 1 & 6 & Cost; the meaning of a resolution \\
\bottomrule
\end{tabularx}
\end{table}

Table~\ref{tab:regime} sets out what the scenario does to the ratios the models depend on. The air over
each square metre, 7.25 times Earth's, stretches every absorber path: water vapour reaches
170~kg\,m$^{-2}$ in the design case and 348~kg\,m$^{-2}$ in its warmest test, where Earth's column holds
about 25 \autocite{trenberth2005trends}.\claim{S08} The scale height, about 52~km against a radius of
1,737~km, makes a poor picture of the Moon's air out of the geometry the radiation codes assume, a stack of
flat layers lit only while the Sun stands above each column's horizon. The air's mass sits on average
48.6~km up, a shell 50~km up has 5.8\% more area than the ground beneath it, and straight rays through the
design column give a horizon air mass of 7.1 for well-mixed air and 13.5 for water vapour, where PlaSim's
formula, written for Earth, gives 35.\claim{S09} The same height that curves the radiation's paths also
strains the dynamics' thin shell: gravity at the 86~km tropopause is 9\% below its surface value, a column
with gravity falling as $GM/r^2$ holds about 5\% more air per unit area than $p_s/g_s$, and the widening
shells hold 11.5\% more over each square metre of ground.\claim{S10}

Gravity and the calendar strain everything else. A sixth of Earth's gravity raises every pressure level
sixfold in height, so a formula written in metres puts its features six times too low; it slows the fall
of droplets sixfold, of rain about threefold and of graupel 2.3-fold; it sets the roughness length of a
wind-driven sea six times longer for the same friction velocity \autocite{charnock1955wind} and slows a
deep-water wave of a given length to 0.41 of its speed on Earth.\claim{S11} The dry dynamics stretch about
sixfold in size and duration at the same speeds, so a 6~km column of a cloud-resolving model stands for
about 1~km at Earth's scale; the microphysics and the radiation keep their own scales, and the gravity pair
meant to test the stretch failed to convect.\claim{S12} The 29.53-day solar day deepens the soil's thermal
wave to about 0.6~m, $\sqrt{29.53}\approx5.4$ times the depth of Earth's diurnal wave, and multiplies nearly
thirtyfold whatever a formula counts per day; the 27.32-day rotation, slow against the Sun but synchronous
with Earth, sets a Coriolis parameter of $\SI{5.2e-6}{s^{-1}}$ at 80\textdegree, an inertial period of
14 Earth days.\claim{S13} The tilt of 1.54\textdegree\ leaves the poles a grazing Sun all year, a regime
that ExoPlaSim's orbit default would replace with Earth's seasons.

The questions each model owns follow from these strains. How warm and humid the air near the ground, where
people and plants would live, runs through the month-long day, how much rain falls and where, and what the
Moon's mean climate becomes under each change of shield or gas: the GCM answers these for the whole Moon,
with clouds diagnosed from Earth formulas. How storms build, how much of their rain reaches the ground and
how the lowest kilometres differ from the GCM's: CM1 answers these where it forms its own clouds. How often
lightning strikes and what each flash carries: the electrified CM1. How high the seas run and what reaches
each shore: SWAN and SWASH, driven by the GCM's stress. The line-by-line model stands behind the
atmospheric models, holding the radiation the GCM is calibrated to and CM1's is checked against.
